\chapter{릴레이와 접점 논리}\label{ch:relay}

1--4부는 프로그램과 디지털 회로의 코드였다. 5부는 전시의 두 번째 기계, 릴레이식 봄베의 회로를 다룬다.
릴레이식 봄베에는 FPGA도, 트랜지스터로 된 논리 게이트도 거의 없다. 계산을 하는 것은 전자석이 끌어당기는
금속 접점이다. 이 장은 릴레이 한 개가 어떻게 생겼는지에서 출발해, 접점을 엮어 논리를 만드는 법, 그리고
그 방법으로 지은 네 가지 회로(위치 표시의 접점 나무, 이진 시프터, 오르기만 하는 계수기)를 차례로 본다.
전기를 몰라도 읽을 수 있게 썼다. 필요한 것은 ``전류는 끊기지 않은 길로만 흐른다''는 사실 하나다.

\section{릴레이 한 개: 코일과 접점}\label{sec:relay-one}

\Why \term{릴레이}{relay}는 전기로 움직이는 스위치다. 철심에 감은 전선(\term{코일}{coil})에 전류가 흐르면
철심이 자석이 되어 얇은 철판(\term{접극자}{armature})을 끌어당긴다. 접극자에는 금속 접점이 붙어 있어서,
끌려가면 접점이 한쪽 단자에서 떨어져 다른 쪽 단자에 닿는다. 전류를 끊으면 스프링이 접극자를 제자리로
돌려보낸다(그림~\ref{fig:r-anatomy}). 코일 쪽 회로와 접점 쪽 회로는 전기적으로 이어져 있지 않다. 작은 코일
전류 하나로 전혀 다른 회로의 길을 열고 닫을 수 있다는 것이 릴레이의 쓸모다.

접점에는 이름이 셋 있다. 움직이는 쪽을 \term{공통}{COM}이라 하고, 코일이 꺼져 있을 때 공통과 닿아 있는
단자를 \term{상시 닫힘}{NC, normally closed}, 코일이 켜져야 닿는 단자를 \term{상시 열림}{NO, normally open}이라
한다. 우리말 문헌에서는 NO를 a접점, NC를 b접점이라고도 부른다. 공통 하나에 NC와 NO가 모두 있는 접점을
\term{전환 접점}{changeover, c접점}이라 한다. 코일 하나가 전환 접점을 몇 벌 움직이는지를 \term{극}{pole}의
수로 센다. 이 프로젝트의 릴레이는 거의 전부 극이 둘인 전환 접점 릴레이, 곧 \term{DPDT}{double pole, double
throw}다.

\begin{figure}[ht]
\centering
\begin{tikzpicture}[x=1cm,y=1cm]
  \foreach \dx/\on/\ttl in {0/0/{코일 꺼짐},7/1/{코일 켜짐}} {
    \begin{scope}[xshift=\dx cm]
      \node[font=\small\bfseries] at (2.2,3.3) {\ttl};
      % core and coil
      \fill[cgray!40] (0,0) rectangle (0.6,2.0);
      \foreach \y in {0.25,0.5,...,1.75} \draw[corange,line width=1.2pt] (-0.1,\y) -- (0.7,\y);
      \node[note,anchor=north] at (0.3,-0.1) {코일};
      \ifnum\on=1
        \foreach \y in {0.5,1.0,1.5} \draw[->,cblue] (0.9,\y) -- (1.25,\y);
      \fi
      % armature (pivot at (0.3,2.3))
      \fill (0.3,2.35) circle (0.06);
      \ifnum\on=0
        \draw[line width=2pt,cgray] (0.3,2.35) -- (3.2,2.75);
        \draw[line width=1pt] (3.2,2.75) -- (3.2,2.35);
        \fill[cblue] (3.2,2.35) circle (0.07);
      \else
        \draw[line width=2pt,cgray] (0.3,2.35) -- (3.2,2.05);
        \draw[line width=1pt] (3.2,2.05) -- (3.2,1.45);
        \fill[cblue] (3.2,1.45) circle (0.07);
      \fi
      % NC (top) and NO (bottom) terminals
      \draw[line width=1pt] (3.0,2.42) -- (4.2,2.42) node[right,font=\footnotesize] {NC};
      \draw[line width=1pt] (3.0,1.38) -- (4.2,1.38) node[right,font=\footnotesize] {NO};
      \draw[line width=1pt] (0.3,2.35) -- (0.3,2.9) -- (-0.4,2.9) node[left,font=\footnotesize] {COM};
      \draw[cgreen] (3.9,2.6) -- (3.8,2.68) -- (4.0,2.76) -- (3.8,2.84) -- (4.0,2.92) -- (3.8,3.0) -- (3.9,3.08);
      \node[note,anchor=west] at (4.0,3.0) {스프링};
      \ifnum\on=0
        \node[note,text=cblue,anchor=west] at (0.9,0.7) {COM--NC가\\닿아 있다};
      \else
        \node[note,text=cblue,anchor=west] at (1.35,0.7) {COM--NO가\\닿아 있다};
      \fi
    \end{scope}
  }
\end{tikzpicture}
\caption{릴레이의 동작. 코일에 전류가 흐르면 철심이 접극자를 끌어당겨 공통(COM)이 NC에서 떨어지고 NO에 닿는다.
코일 전류를 끊으면 스프링이 되돌린다.}
\label{fig:r-anatomy}
\end{figure}

릴레이는 느리다. 코일에 전압을 건 뒤 접점이 실제로 넘어가 흔들림(\term{튐}{bounce})까지 가라앉는 데 이
프로젝트의 소형 신호 릴레이로 2.5--5~ms가 걸리고, 놓일 때도 몇 ms가 걸린다. FPGA의 플립플롭이 수 나노초에
바뀌는 것과 비교하면 백만 배쯤 느리다. 이 느림과, 같은 품번이라도 개체마다 시간이 다르다는 점이
\ref{ch:relay-timing}장의 주제다.

\Ours 릴레이식 봄베의 모든 보드는 Omron G6S-2급의 24~V DPDT 신호 릴레이 한 종류를 쓴다. 회로도 생성기는
핀 번호와 기능의 대응을 한 줄로 적어 두고 모든 보드가 이 줄을 공유한다.

\begin{lstlisting}[language=Python,caption={\code{gen\_shifter.py}: G6S-2의 핀 기능. 코일은 1(+)과 12($-$), 극 1은 4/3/5, 극 2는 9/10/8.},label={lst:r-pins}]
# G6S-2 pin functions (verified against the library symbol graphics: the armature of pole 1
# rests on pin 3, of pole 2 on pin 10, when the coil is off).
RELAY_PINS = dict(coil_p="1", coil_n="12", c1="4", nc1="3", no1="5", c2="9", nc2="10", no2="8")
\end{lstlisting}

검증기도 같은 대응을 따로 적는다(\code{POLES = [("4", "3", "5"), ("9", "10", "8")]}). 생성기와 검증기가 같은
사실을 각자 적어 두는 이유는 \ref{ch:boardverify}장에서 본다.

\section{접점으로 논리를 만든다}\label{sec:relay-logic}

\Why \ref{ch:hdl}장의 게이트는 0과 1을 전압으로 나타냈다. 릴레이 회로에서도 같다. 어떤 선에 +24~V가 닿아
있으면 1, 아니면 0이다. 다만 계산하는 방법이 다르다. 게이트는 입력 전압을 보고 출력 전압을 새로 만들지만,
접점은 길을 열고 닫을 뿐이다. 그래서 접점 회로의 논리는 ``전원에서 출발해 닫힌 접점만 밟아 이 점에 올 수
있는가''라는 길 찾기 문제가 된다. 이 관점은 \ref{ch:boardverify}장의 시뮬레이터가 그대로 쓴다.

\begin{itemize}
\item 두 접점을 \emph{직렬}로 이으면 둘 다 닫혀야 길이 생긴다. AND다.
\item 두 접점을 \emph{병렬}로 이으면 하나만 닫혀도 길이 생긴다. OR다.
\item NC 접점은 코일이 꺼져 있을 때 닫힌다. NOT이다.
\end{itemize}

\General 계단 조명의 3로 스위치, 자동차의 경적 릴레이, 엘리베이터의 층 선택 회로는 모두 이 세 규칙으로
만든 것이다. 그림~\ref{fig:r-logic}는 램프 하나를 세 가지 방식으로 켜는 회로다.

\begin{figure}[ht]
\centering
\begin{circuitikz}[x=1cm,y=1cm]
  \ctikzset{bipoles/length=9mm}
  \foreach \dx/\ttl in {0/AND (직렬),5/OR (병렬),10/NOT (NC 접점)} \node[font=\small\bfseries] at (\dx+1.7,2.4) {\ttl};
  % AND
  \draw (0,1.6) node[left,font=\footnotesize] {$+24$ V} to[nos,l={\scriptsize $a$}] (1.4,1.6) to[nos,l={\scriptsize $b$}] (2.8,1.6) -- (3.3,1.6) to[lamp] (3.3,0) -- (0,0) node[left,font=\footnotesize] {0 V};
  % OR
  \begin{scope}[xshift=5cm]
  \draw (0,1.6) node[left,font=\footnotesize] {$+24$ V} -- (0.4,1.6) coordinate (s);
  \draw (s) to[nos,l={\scriptsize $a$}] (2.2,1.6) coordinate (e);
  \draw (s) -- (0.4,0.9) to[nos,l_={\scriptsize $b$}] (2.2,0.9) -- (e);
  \draw (e) -- (3.3,1.6) to[lamp] (3.3,0) -- (0,0) node[left,font=\footnotesize] {0 V};
  \end{scope}
  % NOT
  \begin{scope}[xshift=10cm]
  \draw (0,1.6) node[left,font=\footnotesize] {$+24$ V} to[ncs,l={\scriptsize $a$}] (2.2,1.6) -- (3.3,1.6) to[lamp] (3.3,0) -- (0,0) node[left,font=\footnotesize] {0 V};
  \end{scope}
\end{circuitikz}
\caption{접점 논리. 접점 $a$, $b$는 각각 다른 릴레이의 접점이다(소문자는 접점, 그 릴레이의 코일은 대문자로
적는다). 열린 스위치 기호가 NO, 닫힌 기호가 NC다.}
\label{fig:r-logic}
\end{figure}

하나 더 있다. \term{다이오드}{diode}는 전류를 한 방향으로만 흘리는 부품이다. 여러 신호를 다이오드 하나씩을
거쳐 한 점에 모으면, 어느 하나라도 +24~V이면 그 점이 +24~V가 되고 거꾸로 그 점의 전압이 다른 입력으로
새어 나가지 않는다. 이것을 \term{다이오드 OR}{diode OR}라고 한다. 접점 하나를 쓰지 않고 OR를 만들 수 있어
이 프로젝트의 계수기는 다이오드를 자리마다 64개나 쓴다.

\Ours 접점 논리가 게이트 논리와 가장 다른 점은 \emph{양방향}이라는 것이다. 게이트는 입력에서 출력으로만
신호를 보내지만, 닫힌 접점은 그냥 전선이라 양쪽으로 전류가 흐른다. 그래서 의도하지 않은 우회로(\term{샛길}{sneak
path})가 생기기 쉽다. 이 프로젝트의 시프터 보드는 ``한 단의 릴레이가 모두 함께 움직인다''는 성질로
샛길을 막고(\ref{sec:relay-shifter}절), 테스트 레지스터는 선마다 다이오드를 넣어 막는다(\ref{sec:t-register}절).
검증기가 ``기대한 연결 \emph{말고는} 아무것도 이어지지 않았는가''를 꼭 확인하는 이유도 이것이다.

\section{자기 유지: 릴레이로 만든 기억}\label{sec:relay-latch}

\Why 릴레이의 접점이 자기 코일의 전원 길에 들어가면 기억이 생긴다. 한 번 켜진 릴레이가 자기 NO 접점으로
자기 코일에 전류를 계속 대 주기 때문이다. 처음 켠 신호가 사라져도 릴레이는 켜진 채 남는다. 끄려면 그
자기 유지 길 어딘가를 따로 끊어야 한다. 이것이 \term{자기 유지}{self-hold} 회로이고, \ref{ch:sync}장의
플립플롭이 하던 ``1비트 기억''을 릴레이 하나와 접점 하나로 만든 것이다.

\General 공장 기계의 시작/정지 버튼이 전형적인 예다(그림~\ref{fig:r-latch}). 시작 버튼(NO)을 누르면
릴레이 $K$가 켜지고, $K$의 접점 $k$가 시작 버튼과 병렬로 닫혀 버튼을 놓아도 $K$가 남는다. 정지 버튼(NC)은
이 길 전체와 직렬이어서, 누르면 길이 끊겨 $K$가 놓이고 $k$도 열린다. 정전 뒤에 기계가 저절로 다시 돌지
않는다는 안전상의 이점도 있다.

\begin{figure}[ht]
\centering
\begin{circuitikz}[x=1cm,y=1cm]
  \ctikzset{bipoles/length=9mm}
  \draw (0,1.6) node[left,font=\footnotesize] {$+24$ V} to[ncs,l={\scriptsize 정지}] (1.6,1.6) coordinate (a)
        to[nos,l={\scriptsize 시작}] (3.4,1.6) coordinate (b) -- (4.2,1.6);
  \draw (a) -- (1.6,0.7) to[nos,l_={\scriptsize $k$ (자기 유지)}] (3.4,0.7) -- (b);
  \node[rcoil,anchor=west] (K) at (4.2,1.6) {$K$};
  \draw (K.east) -- ++(0.4,0) node[right,font=\footnotesize] {0 V};
  \node[circ] at (a) {}; \node[circ] at (b) {};
\end{circuitikz}
\hspace{6mm}
\begin{tikztimingtable}[timing/slope=0.1,timing/coldist=4pt,xscale=1.5,yscale=1.2,semithick]
  시작 & 1L .6H 4.4L \\
  정지 & 4L .6H 1.4L \\
  $K$   & 1.2L 2.9H 1.9L \\
\end{tikztimingtable}
\caption{시작/정지 회로. 시작 버튼을 잠깐 눌러도 $K$는 자기 접점 $k$로 켜진 채 남고, 정지 버튼이 길을 끊어야
놓인다. 오른쪽의 $K$가 버튼보다 조금 늦게 바뀌는 것은 릴레이의 동작 시간이다.}
\label{fig:r-latch}
\end{figure}

\Ours 자기 유지는 이 프로젝트의 릴레이 회로 곳곳에 있다.
\begin{itemize}
\item 테스트 레지스터의 감지 릴레이 26개는 시험 전압이 닿은 짧은 순간에 붙고, 자기 유지로 그 결과를 붙든
뒤 판정을 받는다(\ref{sec:t-register}절). 유지 길을 공통의 \term{유지 버스}{hold bus}에서 먹이므로 그 버스
하나를 끊으면 26개가 한꺼번에 지워진다.
\item 계수기의 슬레이브 릴레이는 자기 유지로 값을 기억한다(\ref{sec:relay-counter}절).
\item 정지 릴레이는 정지를 붙들어 기계를 세워 두고, 해제 버튼이 그 유지를 끊는다.
\end{itemize}

\section{먼저 끊고 나중에 잇기}\label{sec:relay-bbm}

\Why 전환 접점의 공통은 NC에서 떨어져 NO로 \emph{날아간다}. 그 사이 잠깐(이 프로젝트의 시뮬레이터는
0.2--0.6~ms로 잡는다) 공통은 어디에도 닿아 있지 않다. 이 성질을 \term{먼저 끊고 나중에 잇기}{break-before-make}라
한다. 덕분에 NC 쪽 회로와 NO 쪽 회로는 \emph{결코 동시에} 이어지지 않는다. 한 접점이 두 길 가운데 정확히
하나를 고르는 것이 기계적으로 보장되는 것이다.

릴레이 두 개로 같은 일을 하면 이 보장이 사라진다. 하나는 켜지고 하나는 꺼지는 순간에 두 릴레이의 시간이
다르면, 잠깐 둘 다 닫혀 있거나 둘 다 열려 있게 된다(그림~\ref{fig:r-bbm}).

\begin{figure}[ht]
\centering
\begin{tikztimingtable}[timing/slope=0.1,timing/coldist=4pt,xscale=1.6,yscale=1.2,semithick]
  코일 & 1L 4H 1L \\
  NC 쪽 & 1.6H 4L .4H \\
  NO 쪽 & 1.8L 3.4H .8L \\
  \textcolor{cred}{릴레이 둘: A} & 1.4H 4L .6H \\
  \textcolor{cred}{릴레이 둘: B} & 1.2L 4.6H .2L \\
\extracode
  \begin{pgfonlayer}{background}
    \fill[cgreenbg] (1.6,-5.0) rectangle (1.8,1.0);
    \fill[credbg] (1.2,-9.6) rectangle (1.4,-5.4);
    \fill[credbg] (5.4,-9.6) rectangle (5.8,-5.4);
  \end{pgfonlayer}
  \node[note,text=cgreen,anchor=west] at (6.2,-2.2) {전환 접점 하나: 끊긴 틈(초록)은\\있어도 겹침은 없다};
  \node[note,text=cred,anchor=west] at (6.2,-7.6) {따로 움직이는 두 릴레이:\\시간이 어긋나면 둘 다 닫힌다(빨강)};
\end{tikztimingtable}
\caption{먼저 끊고 나중에 잇기. 위는 전환 접점 하나의 NC와 NO, 아래는 같은 일을 릴레이 두 개(A는 NC처럼,
B는 NO처럼 쓴 것)에 맡긴 경우다.}
\label{fig:r-bbm}
\end{figure}

\General 철도의 선로 전환기가 같은 원리다. 레일이 한쪽에서 떨어진 뒤에야 다른 쪽에 붙으므로 열차가 두 선로에
동시에 이어지는 일이 없다. 전원 회로에서는 정반대 성질(먼저 잇고 나중에 끊기)이 필요할 때도 있다. 그래서
부품 자료에는 어느 쪽인지가 꼭 적혀 있다.

\Ours 릴레이식 계수기는 클럭을 \code{CK}와 \code{NCK} 두 선으로 받는다. 둘은 반드시 정반대여야 하고, 둘이
동시에 켜지면 계수기 값이 틀어진다. 그래서 순서기는 이 두 선을 \emph{릴레이 하나의 한 극}에서 낸다. 공통에
+24~V, NO가 \code{CK}, NC가 \code{NCK}다. 계수기 보드의 설명서는 ``개별 릴레이 둘로 만들어 둘 다 켜지면 값이
틀어진다(시뮬레이션에서 확인)''고 적고 있고, 순서기의 구조 검사는 \code{CK}와 \code{NCK}가 버스마다 한
릴레이의 한 극에서만 나오는지를 확인한다. 끊긴 틈 동안에는 둘 다 꺼지는데, 이 틈이 릴레이가 놓이는 시간보다
훨씬 짧아 계수기의 릴레이들은 그동안 상태를 유지한다.

\section{역기전력 다이오드: 코일을 끌 때}\label{sec:relay-flyback}

\Why 코일은 전류의 변화를 싫어한다. 코일에 흐르던 전류를 갑자기 끊으면 코일은 그 전류를 계속 흘리려고
매우 큰 전압을 만든다(\term{역기전력}{back EMF}). 전압의 크기는 전류가 줄어드는 빠르기에 비례하므로,
스위치가 순식간에 열리면 수백 V가 나올 수도 있다. 이 전압은 코일을 끈 스위치(접점이나 트랜지스터)를
상하게 하고 주변 회로에 잡음을 뿌린다.

해법은 코일 양 끝에 다이오드를 거꾸로 거는 것이다(그림~\ref{fig:r-flyback}). 평소에는 다이오드가 막혀 있어
아무 일도 하지 않는다. 스위치가 열리는 순간 코일의 전류는 다이오드를 지나 코일로 되돌아가는 작은 고리를
돌며 천천히 사라진다. 코일 양 끝의 전압은 다이오드 하나의 전압(약 0.7~V)으로 묶인다. 이 다이오드를
\term{역기전력 다이오드}{flyback diode}라 한다.

\begin{figure}[ht]
\centering
\begin{circuitikz}[x=1cm,y=1cm]
  \ctikzset{bipoles/length=9mm}
  \foreach \dx/\ttl in {0/켜져 있을 때,6.5/스위치를 연 직후} {
  \begin{scope}[xshift=\dx cm]
    \node[font=\small\bfseries] at (1.9,3.3) {\ttl};
    \draw (0,2.5) node[left,font=\footnotesize] {$+24$ V} -- (0.3,2.5);
    \ifdim\dx pt=0pt
      \draw (0.3,2.5) to[cspst] (1.6,2.5);
    \else
      \draw (0.3,2.5) to[ospst] (1.6,2.5);
    \fi
    \draw (1.6,2.5) -- (2.6,2.5) to[L,l_={코일}] (2.6,0) -- (0,0) node[left,font=\footnotesize] {0 V};
    \draw (2.6,2.5) -- (3.8,2.5);
    \draw (3.8,0) to[D*,l_={\scriptsize 다이오드}] (3.8,2.5);
    \draw (3.8,0) -- (2.6,0);
    \node[circ] at (2.6,2.5) {}; \node[circ] at (2.6,0) {};
  \end{scope}}
  \draw[->,cblue,line width=1pt] (1.2,2.8) -- (2.2,2.8);
  \draw[->,cblue,line width=1pt] (2.35,2.1) -- (2.35,0.6);
  \draw[->,corange,line width=1pt] (8.85,2.1) -- (8.85,0.5);
  \draw[->,corange,line width=1pt] (9.4,-0.25) -- (10.0,-0.25);
  \draw[->,corange,line width=1pt] (10.55,0.5) -- (10.55,2.0);
  \draw[->,corange,line width=1pt] (10.0,2.75) -- (9.4,2.75);
  \node[note,text=corange,anchor=west] at (11.8,1.3) {코일의 전류가\\다이오드로 돌며\\사라진다};
\end{circuitikz}
\caption{역기전력 다이오드. 다이오드는 평소에 막혀 있다(삼각형이 가리키는 쪽이 $+$). 스위치가 열리면 코일의
전류가 다이오드를 지나는 고리로 돌아 천천히 줄어든다.}
\label{fig:r-flyback}
\end{figure}

대가가 있다. 전류가 천천히 사라지므로 릴레이가 \emph{늦게 놓인다}. 이 프로젝트는 다이오드만 단 코일의
복귀 시간을 데이터시트 값의 2--3배인 3--10~ms로 추정한다. 빨리 놓여야 하는 릴레이에는 다이오드에 제너
다이오드(정해진 전압을 넘으면 거꾸로도 흐르는 다이오드)를 직렬로 더한다. 코일 전압이 0.7~V가 아니라
약 31~V에서 묶이므로 전류가 훨씬 빨리 줄어 복귀가 1.5--4~ms로 짧아진다. 에너지를 빨리 버리는 대신 순간
전압이 커지는 맞바꿈이다.

\General 트랜지스터로 릴레이나 모터를 켜고 끄는 회로라면 어디에나 이 다이오드가 있다. 아두이노로 릴레이를
움직이는 입문 예제에서 다이오드를 빠뜨리면 트랜지스터가 몇 번 만에 망가진다.

\Ours 이 프로젝트의 모든 릴레이 코일에는 역기전력 다이오드(1N4148W)가 하나씩 있다. 시프터 보드 한 장에
릴레이 130개와 다이오드 130개가 있는 이유다. 생성기는 릴레이를 놓을 때마다 다이오드를 같이 놓는다.

\begin{lstlisting}[language=Python,caption={\code{gen\_shifter.py}: 릴레이마다 코일 양 끝에 다이오드(음극이 $+$ 쪽).},label={lst:r-flyback}]
        sh.pin_net(S_RELAY, x, y, P["coil_p"], "CTRL")
        sh.pin_net(S_RELAY, x, y, P["coil_n"], "GND")
        # flyback diode across the coil (cathode to the + side)
        dref = f"D{st}{i + 1:02d}"
        dx = g(x - 24.0)
        sh.symbol(S_DIODE, dx, y, dref, "1N4148W", FP_DIODE, "", "Coil flyback diode", path, show_value=False)
        sh.pin_net(S_DIODE, dx, y, "1", "CTRL")   # K
        sh.pin_net(S_DIODE, dx, y, "2", "GND")    # A
\end{lstlisting}

테스트 레지스터의 타이밍 릴레이 셋(\code{STROBE}, \code{JUDGE}, \code{CLEAR})은 빨리 놓여야 하므로 다이오드 +
제너 30~V를 쓴다. 그러면 새 문제가 생긴다. 신호를 끄는 순간 코일의 $+$ 쪽이 약 $-31$~V까지 내려가는데, 그
$+$ 쪽이 곧 커넥터의 핀이면 그 음전압이 케이블을 타고 신호를 내 준 쪽(기어식의 광센서 출력 등)으로
나간다. 그래서 결정 0003은 ``받는 쪽이 역기전력을 자기 안에서 처리한다''는 규칙을 두고, 핀과 코일 사이에
다이오드를 하나 더 넣었다. 끄는 순간 이 다이오드가 막혀 코일의 전류는 보드 안의 클램프로만 돈다. 코일
전압은 다이오드 강하 0.7~V만큼 줄어든다.

\section{접점 나무: 5비트를 26개 가운데 하나로}\label{sec:relay-tree}

릴레이 계수기는 위치를 5비트 이진수로 갖는다. 관람객에게는 그것을 글자 하나로 보여 주어야 한다. 5개의
선(비트)을 보고 26개 램프 가운데 정확히 하나를 켜는 회로를 \term{해독기}{decoder}라 한다.

\Why 전환 접점 하나는 길 하나를 둘로 가른다. 그러면 비트마다 전환 접점을 한 단씩 두고 이어 나가면
가지가 갈라지는 나무가 된다. 뿌리에 +24~V를 걸면, 각 단에서 그 비트의 값에 따라 한 갈래만 닫혀 있으므로
전압은 정확히 한 잎에만 닿는다(그림~\ref{fig:r-tree}). 갈림마다 갈래가 정확히 하나라는 것은 전환 접점의
성질(\ref{sec:relay-bbm}절)이어서, 릴레이들이 서로 다른 순간에 넘어가는 전환 중에도 ``아무것도 안 켜짐''이나
``둘이 켜짐''이 생기지 않는다. 잠깐 다른 글자가 켜질 수는 있다.

\begin{figure}[ht]
\centering
\begin{tikzpicture}[x=1cm,y=0.62cm,
    nd/.style={circle,fill=black,inner sep=1.2pt},
    lf/.style={draw,rounded corners=1pt,font=\scriptsize\ttfamily,inner sep=2pt,minimum width=12mm},
    on/.style={line width=1.6pt,cblue},
    br/.style={font=\tiny,text=cgray,fill=white,inner sep=0.5pt}]
  \node[font=\footnotesize,anchor=east] (root) at (0,0) {$+24$ V};
  \coordinate (r) at (0.1,0);
  % level 1 nodes (after bit 0): p=0 (NC), p=1 (NO)
  \coordinate (n0) at (2.6,3.0); \coordinate (n1) at (2.6,-3.0);
  % level 2 nodes (after bit 1): from n0: p=0, p=2 ; from n1: p=1, p=3
  \coordinate (m0) at (5.2,4.5); \coordinate (m2) at (5.2,1.5);
  \coordinate (m1) at (5.2,-1.5); \coordinate (m3) at (5.2,-4.5);
  % leaves (after bit 2): from mp: p (NC), p+4 (NO)
  \foreach \p/\y in {0/5.25,4/3.75,2/2.25,6/0.75,1/-0.75,5/-2.25,3/-3.75,7/-5.25} {
    \node[lf,anchor=west] (L\p) at (7.8,\y) {램프 \p};
  }
  \draw (r) -- node[br,pos=0.55] {NC} (n0); \draw[on] (r) -- node[br,pos=0.55] {NO} (n1);
  \draw (n0) -- node[br] {NC} (m0); \draw (n0) -- node[br] {NO} (m2);
  \draw[on] (n1) -- node[br] {NC} (m1); \draw (n1) -- node[br] {NO} (m3);
  \foreach \p/\hi in {0/4,2/6,1/5,3/7} {
    \draw (m\p) -- node[br,pos=0.45] {NC} ([xshift=-8mm]L\p.west) -- (L\p.west);
    \draw (m\p) -- node[br,pos=0.45] {NO} ([xshift=-8mm]L\hi.west) -- (L\hi.west);
  }
  \draw[on] (m1) -- ([xshift=-8mm]L5.west) -- (L5.west);
  \node[lf,anchor=west,fill=cyellowbg,draw=corange] at (L5.west) {램프 5};
  \foreach \c in {r,n0,n1,m0,m1,m2,m3} \node[nd] at (\c) {};
  \node[font=\footnotesize\bfseries] at (1.3,6.4) {비트 0};
  \node[font=\footnotesize\bfseries] at (3.9,6.4) {비트 1};
  \node[font=\footnotesize\bfseries] at (6.4,6.4) {비트 2};
  \node[note,anchor=west,align=left,text=cblue] at (9.4,-2.25) {값 5 = 101$_2$:\\비트 0 = 1, 비트 1 = 0, 비트 2 = 1};
\end{tikzpicture}
\caption{3비트 접점 나무(잎 8개). 뿌리를 가장 낮은 비트에 두었다. 마디마다 전환 접점 하나가 있고, 그 비트가
1이면 NO 갈래, 0이면 NC 갈래가 닫힌다. 파란 길이 값 5의 길이다. 잎의 순서가 0, 4, 2, 6, \ldots로 섞이는 것은
가장 낮은 비트부터 갈랐기 때문이다.}
\label{fig:r-tree}
\end{figure}

\General 옛 전화 교환기의 선택 회로, 그리고 디지털 회로의 \term{멀티플렉서 나무}{multiplexer tree}가 같은
모양이다. 잎이 $n$개인 두 갈래 나무에는 갈림이 정확히 $n-1$개 있다. 갈림 하나에 전환 접점 하나이므로
26글자에는 접점 25개가 든다.

\Ours 위치 표시 보드(\code{hardware/relay/position-display})는 이 나무를 세 자리($l$, $m$, $r$)에 하나씩 둔다.
뿌리를 어느 비트에 두느냐가 릴레이 수를 정한다. 한 비트의 접점은 모두 그 비트의 코일로 움직여야 하는데,
G6S-2 하나에는 극이 둘뿐이므로 접점이 많은 비트는 릴레이 여러 개가 코일을 나란히 달고 함께 움직인다.

\begin{center}
\small
\begin{tabular}{lccccccc}
\toprule
뿌리 & 비트 0 & 비트 1 & 비트 2 & 비트 3 & 비트 4 & 접점 & 2극 릴레이 \\
\midrule
가장 높은 비트 & 13 & 6 & 3 & 2 & 1 & 25 & 7+3+2+1+1 = 14 \\
가장 낮은 비트 & 1 & 2 & 4 & 8 & 10 & 25 & 1+1+2+4+5 = \textbf{13} \\
\bottomrule
\end{tabular}
\end{center}

가장 낮은 비트를 뿌리로 하면 단 $b$에는 아래 $b$개 비트의 값 $p$마다 접점이 하나여서 $2^b$개다. 마지막 단(비트
4, 값 16)은 $16+p \le 25$인 $p = 0..9$에만 접점이 필요하다. $p = 10..15$의 마디는 그대로 램프 10--15가 된다.
생성기는 이 규칙을 마디 이름을 짓는 함수 하나로 적는다.

\begin{lstlisting}[language=Python,caption={\code{gen\_display.py}: 나무의 마디와 접점. 마디 이름에 ``지금까지 지나온 비트의 값'' $p$가 들어 있다.},label={lst:r-tree}]
CONTACTS = [1, 2, 4, 8, 10]
RELAYS = [(c + 1) // 2 for c in CONTACTS]        # G6S-2 = 2 poles -> 1, 1, 2, 4, 5 = 13 per digit

def node(D: str, b: int, p: int) -> str:
    """Tree node at level b (relay contacts of bit b have their COM here) with prefix p = value of bits 0..b-1."""
    if b == 0:
        return "+24V"
    if b == 4 and p >= 10:          # bit 4 has no contact here: 16 + p > 25, the node is the lamp leaf itself
        return leaf(D, p)
    return f"{D}N{b}_{p:02d}"
...
        for b in range(BITS):
            prefixes = list(range(2 ** b)) if b < 4 else list(range(10))
            for p in prefixes:
                com = node(D, b, p)
                if b < 4:
                    nc, no = node(D, b + 1, p), node(D, b + 1, p + 2 ** b)
                else:
                    nc, no = leaf(D, p), leaf(D, p + 16)
\end{lstlisting}

읽는 법은 이렇다. 단 $b$의 접점은 공통이 마디 $(b, p)$에 있고, 그 비트가 0이면(NC) 다음 단의 같은 $p$로,
1이면(NO) $p + 2^b$로 간다. 다섯 단을 지나면 $p$는 비트를 모두 더한 값, 곧 위치의 값이 된다. 계수기는 0--25만
내므로 값 26--31은 생기지 않는다. 생긴다면 접점이 없는 마디로 가서 램프 K--P가 켜진다. 이 보드는 자리마다
릴레이 13개, 세 자리 39개를 쓰고, 검증기는 26$^3$ = 17,576가지 위치 모두에서 정확히 램프 세 개만 켜짐을
확인한다(\ref{ch:boardverify}장).

\section{이진 시프터: 회전을 배선으로}\label{sec:relay-shifter}

이제 봄베의 본체인 스크램블러다. 기어식에서는 드럼이 실제로 돌지만 릴레이식에는 도는 것이 없다. 로터의
회전을 접점으로 흉내 내야 한다.

\Why 로터는 고정된 배선 $W$가 든 원판이다. 원판이 $p$칸 돌아가 있으면, 들어온 글자는 먼저 $p$칸 밀린 자리로
배선에 들어가고, 나올 때 다시 $p$칸 되돌려진다. ``$p$칸 밀기''를 $S(p)$라 쓰면 위치 $p$의 로터는
\[
  S(-p)\;W\;S(p) \qquad\text{(오른쪽부터 읽는다: 밀고, 배선을 지나고, 되돌린다)}
\]
다. \ref{sec:through}절의 \code{through} 함수가 ``들어갈 때 $s$를 더하고, 표를 읽고, 나올 때 $s$를 뺀다''로
하던 일을 그대로 옮긴 것이다.
그러면 회전하는 로터는 \emph{움직이지 않는 배선 카드 하나}와 \emph{밀기 장치 둘}로 바꿀 수 있다. 배선은
인쇄 회로 기판의 동선이면 되고, 밀기만 릴레이로 만들면 된다.

밀기는 어떻게 만드는가. 26가지 밀기를 위치마다 따로 배선하면(위치 릴레이 26개 가운데 하나만 동작) 로터 한
단에 접점이 $26 \times 26$개 든다. 대신 이진수를 쓴다. 0--25는 다섯 비트로 쓸 수 있고, ``$p$칸 밀기''는
``1칸, 2칸, 4칸, 8칸, 16칸 밀기 가운데 $p$의 비트가 1인 것들을 차례로 하기''와 같다. 덧셈은 순서를 바꿔도
되므로 단의 순서는 상관없다. 단 하나는 26선 각각에 전환 접점 하나다. 공통이 입력 $i$, NC가 출력 $i$, NO가
출력 $i+s$다(그림~\ref{fig:r-stage}). 한 단의 26개 접점이 모두 같은 코일 신호로 함께 넘어가므로, 어떤 출력에서
두 입력이 만나는 일(샛길)이 없다.

\begin{figure}[ht]
\centering
\begin{tikzpicture}[x=1cm,y=0.55cm]
  \foreach \dx/\on/\ttl in {0/0/{코일 꺼짐: 그대로},7.5/1/{코일 켜짐: $i \to i+2$}} {
    \begin{scope}[xshift=\dx cm]
      \node[font=\small\bfseries] at (2,1.1) {\ttl};
      \foreach \i/\l in {0/A,1/B,2/C,3/D,4/E,5/F,6/G,7/H} {
        \node[font=\scriptsize\ttfamily,anchor=east] at (0,-\i) {\l};
        \fill (0.1,-\i) circle (1.3pt);
        \fill (3.9,-\i) circle (1.3pt);
        \node[font=\scriptsize\ttfamily,anchor=west] at (4.0,-\i) {\l};
      }
      \foreach \i in {0,...,7} {
        \pgfmathtruncatemacro{\j}{\i+2}
        \ifnum\on=0
          \draw[line width=1.1pt,cblue] (0.1,-\i) -- (3.9,-\i);
          \ifnum\j<8 \draw[dashed,cgray] (0.1,-\i) -- (3.9,-\j); \else \draw[dashed,cgray] (0.1,-\i) -- (2.6,-\i-1.4); \fi
        \else
          \draw[cgray!60] (0.1,-\i) -- (3.9,-\i);
          \ifnum\j<8 \draw[line width=1.1pt,cblue] (0.1,-\i) -- (3.9,-\j); \else \draw[line width=1.1pt,cblue] (0.1,-\i) -- (2.6,-\i-1.4); \fi
        \fi
      }
      \node[note,anchor=north] at (2.6,-8.5) {G, H $\to$ I, J (그림 밖). Y, Z $\to$ A, B (mod 26)};
    \end{scope}}
\end{tikzpicture}
\caption{이진 시프터의 한 단(이동량 2, 26선 가운데 8선). 선마다 전환 접점 하나: 공통 = 입력, NC = 같은 글자,
NO = 두 칸 뒤 글자. 파란 선이 닫힌 갈래다. 단 다섯(1, 2, 4, 8, 16칸)을 직렬로 이으면 0--25칸 어떤 밀기도
된다.}
\label{fig:r-stage}
\end{figure}

\General CPU 안에서 비트를 여러 칸 미는 \term{배럴 시프터}{barrel shifter}가 같은 생각이다. 1, 2, 4, \ldots칸
미는 단을 차례로 두고 각 단을 켤지 말지를 미는 양의 비트로 정한다. 단의 수는 미는 양의 비트 수만큼만 든다.

\Ours 기준 모델은 Rust 코어의 \code{relay.rs}에 있다. 다섯 단을 차례로 지나는 함수는 그림 그대로다.

\begin{lstlisting}[language=Rust,caption={\code{relay.rs}: 이진 5단 시프터. \code{reverse}는 모듈을 거꾸로 끼운 경우(반대로 민다).},label={lst:r-shiftbin}]
pub fn shift_binary(c: u8, control: u8, reverse: bool) -> u8 {
    debug_assert!(control < 26);
    let mut x = c;
    for i in 0..5 {
        if control >> i & 1 == 1 {
            let step = 1u8 << i;
            x = if reverse { alphabet::sub(x, step) } else { shift(x, step) };
        }
    }
    x
}
\end{lstlisting}

\paragraph{밀기를 합친다.} 스크램블러는 오른쪽, 가운데, 왼쪽 로터를 지나 반사판에서 돌아 다시 세 로터를
거꾸로 지난다. 신호가 지나는 순서대로 가는 길 반쪽을 펼치면 로터마다 ``밀고, 배선, 되돌림''이 이어지는데,
이웃한 ``되돌림''과 ``밀기''는 하나로 합칠 수 있다. 오른쪽, 가운데, 왼쪽 로터의 위치를 $r$, $m$, $l$, 크립
오프셋을 $k$라 하면(이번에는 신호가 지나는 순서대로 \emph{왼쪽부터} 읽는다)
\[
  S(r{+}k)\,W_R\,S(-(r{+}k))\;S(m)\,W_M\,S(-m)\;S(l)\,W_L\,S(-l)
  \;=\; S(a_1)\,W_R\,S(a_2)\,W_M\,S(a_3)\,W_L\,S(-l)
\]
\[
  a_1 = r+k,\qquad a_2 = m-(r+k),\qquad a_3 = l-m \qquad(\text{모두 mod 26})
\]
이 된다. 시프터가 반쪽에 여섯에서 넷으로 준다. 마지막 $S(-l)$은 ``$l$칸 밀기'' 보드를 거꾸로 끼워 만든다.
접점은 양방향이므로 $i \to i+l$로 배선한 보드를 반대쪽에서 보면 $i \to i-l$이다. 뺄셈 회로가 필요 없다.

\paragraph{오는 길은 같은 보드를 거꾸로.} 반사판에서 돌아오는 길은 가는 길의 역이다. 시프터 보드 한 장에는
가는 길 반쪽과 오는 길 반쪽의 같은 단이 \emph{한 릴레이의 두 극}으로 들어 있다. 극 1은 $i \to i+s$, 극 2는
$i \to i-s$로 배선해 같은 코일로 움직이므로, 제어값 하나가 가는 길의 밀기와 오는 길의 되돌림을 함께
정한다. 생성기는 이것을 여섯 줄로 적는다.

\begin{lstlisting}[language=Python,caption={\code{gen\_shifter.py}: 단(이동량 \code{shift})의 선 $i$를 맡는 릴레이 하나의 배선.},label={lst:r-stage}]
    for i in range(N):
        ...
        ref = f"K{st}{i + 1:02d}"
        P = RELAY_PINS
        sh.pin_net(S_RELAY, x, y, P["c1"], f"OIN_{letter(i)}")
        sh.pin_net(S_RELAY, x, y, P["nc1"], f"OOUT_{letter(i)}")
        sh.pin_net(S_RELAY, x, y, P["no1"], f"OOUT_{letter((i + shift) % N)}")
        sh.pin_net(S_RELAY, x, y, P["c2"], f"RIN_{letter(i)}")
        sh.pin_net(S_RELAY, x, y, P["nc2"], f"ROUT_{letter(i)}")
        sh.pin_net(S_RELAY, x, y, P["no2"], f"ROUT_{letter((i - shift) % N)}")
\end{lstlisting}

\code{O}는 가는 길(outbound), \code{R}는 오는 길(return)이다. 핀마다 ``이 핀은 이 이름의 선에 닿는다''만 적는다.
이름이 같은 선끼리 이어진다는 KiCad의 규칙이 실제 연결을 만든다(\ref{ch:boardgen}장).

유닛 하나는 시프터 보드 넷, 로터 배선 카드 셋, 반사판 카드 하나로 이루어진다(그림~\ref{fig:r-unit}). 카드는
모두 같은 외곽과 같은 커넥터 자리를 갖고 앞뒤로 겹쳐 꽂히며, 이웃 카드의 같은 자리 커넥터끼리 곧은 리본
케이블로 잇는다. 휠 순서를 바꾸는 것은 로터 카드를 바꿔 꽂는 일이다.

\begin{figure}[ht]
\centering
\begin{tikzpicture}[x=1cm,y=1cm,
    sh/.style={blkb,minimum width=13mm,minimum height=8mm,font=\footnotesize},
    wc/.style={blkg,minimum width=11mm,minimum height=8mm,font=\footnotesize}]
  \node[font=\footnotesize,anchor=east] (in) at (0,1) {IN};
  \node[sh] (s1) at (1.0,1) {$S(a_1)$};
  \node[wc] (wr) at (2.5,1) {$W_R$};
  \node[sh] (s2) at (4.0,1) {$S(a_2)$};
  \node[wc] (wm) at (5.5,1) {$W_M$};
  \node[sh] (s3) at (7.0,1) {$S(a_3)$};
  \node[wc] (wl) at (8.5,1) {$W_L$};
  \node[sh] (s4) at (10.0,1) {$S(-l)$};
  \node[sh] (r1) at (1.0,-0.4) {$S(-a_1)$};
  \node[wc] (rr) at (2.5,-0.4) {$W_R^{-1}$};
  \node[sh] (r2) at (4.0,-0.4) {$S(-a_2)$};
  \node[wc] (rm) at (5.5,-0.4) {$W_M^{-1}$};
  \node[sh] (r3) at (7.0,-0.4) {$S(-a_3)$};
  \node[wc] (rl) at (8.5,-0.4) {$W_L^{-1}$};
  \node[sh] (r4) at (10.0,-0.4) {$S(l)$};
  \node[blko,minimum width=12mm,minimum height=22mm,font=\footnotesize,align=center] (u) at (11.7,0.3) {반사판\\$U$};
  \node[font=\footnotesize,anchor=east] (out) at (0,-0.4) {OUT};
  \draw[->] (in) -- (s1); \draw[->] (s1) -- (wr); \draw[->] (wr) -- (s2); \draw[->] (s2) -- (wm);
  \draw[->] (wm) -- (s3); \draw[->] (s3) -- (wl); \draw[->] (wl) -- (s4); \draw[->] (s4) -- (u.west|-s4);
  \draw[->] (u.west|-r4) -- (r4); \draw[->] (r4) -- (rl); \draw[->] (rl) -- (r3); \draw[->] (r3) -- (rm);
  \draw[->] (rm) -- (r2); \draw[->] (r2) -- (rr); \draw[->] (rr) -- (r1); \draw[->] (r1) -- (out);
  \foreach \c in {1,2,3,4} \draw[dashed,cblue] (s\c.south) -- (r\c.north);
  \node[note,anchor=north] at (5.5,-1.0) {파란 점선: 같은 시프터 보드(릴레이 130개)의 두 극. 초록: 수동 배선 카드(같은 카드의 두 벌)};
  \node[note,anchor=south] at (5.5,1.6) {제어: $a_1 = r+k$, $a_2 = m-(r+k)$는 유닛마다, $a_3 = l-m$, $l$은 모든 유닛이 공유};
\end{tikzpicture}
\caption{릴레이 스크램블러 유닛 하나. 시프터 보드 넷이 가는 길과 오는 길을 함께 맡는다. 오른쪽 끝의 보드는
거꾸로 끼워 $-l$칸을 민다.}
\label{fig:r-unit}
\end{figure}

\code{relay.rs}의 시험은 이 분해가 옳은지를 Enigma~I의 휠 순서 120개(반사판 B, C × 로터 순서 60) 모두, 크립
오프셋 세 가지($k = 0, 7, 25$)로 확인한다. 17,576가지 위치 전부에서 릴레이 모델의 치환표가 \ref{ch:rust-bombe}장의
표 조회 스크램블러와 같아야 통과한다. 이 모델이 릴레이 회로의 직접적인 기준이 된다.

접점의 수를 세어 보면 이 설계의 값어치가 보인다. 위치마다 따로 배선하는 방식은 유닛당 4,056개, 이진
시프터를 합치지 않으면 1,560개, 합치면 1,040개(DPDT 릴레이 520개)다.

\section{오르기만 하는 계수기}\label{sec:relay-counter}

시프터는 제어값 $a_1$, $a_2$, $a_3$, $l$을 5비트 이진수로 받는다. 이 값을 매 위치마다 만들어 주는 것이
계수기다. 그런데 $a_2 = m - (r+k)$와 $a_3 = l - m$은 뺄셈이다. 릴레이로 뺄셈기를 만들면 릴레이가 많이 든다.

\Why 뺄셈을 하지 않고 \emph{차이 자체를 센다}. 봄베의 위치는 가장 빠른 $l$이 매 클럭 하나씩 오르고,
$l$이 25에서 0으로 넘어갈 때 $m$이 오르고, $m$이 넘어갈 때 $r$이 오른다. 그러면 차이 $a_3 = l - m$은
\begin{itemize}
\item $l$만 오르는 클럭에는 1 오르고,
\item $l$이 넘어가며 $m$도 오르는 클럭에는 ($+1$과 $-1$이 상쇄되어) 그대로다.
\end{itemize}
$a_2$와 $a_1$도 같은 식으로 ``오르거나 그대로''만 한다. 어느 값도 내려가지 않는다. 그래서 모든 계수기가
\emph{1 더하기 아니면 멈춤}만 하면 되고, 뺄셈기도 오르내리는 계수기도 필요 없다.

\begin{lstlisting}[language=Rust,caption={\code{relay.rs}: 한 클럭. 더하기(mod 26)와 멈춤만 쓴다.},label={lst:r-counters}]
    pub fn clock(&mut self) {
        let l_wraps = self.l == 25;
        let m_advances = l_wraps;
        let m_wraps = m_advances && self.m == 25;
        let c = &mut self.controls;
        c.a4_l = shift(c.a4_l, 1);
        if !l_wraps {
            c.a3 = shift(c.a3, 1);
        }
        if m_advances && !m_wraps {
            c.a2 = shift(c.a2, 1);
        }
        if m_wraps {
            c.a1 = shift(c.a1, 1);
        }
        ...
    }
\end{lstlisting}

시작값은 위치 0에서의 값, 곧 $a_1 = k$, $a_2 = -k$, $a_3 = 0$, $l = 0$이다. 시작값을 넣는 것을 \term{프리셋}{preset}이라
한다. Rust 시험 \code{counters\_track\_controls\_over\_full\_run}은 크립 오프셋 26가지 모두에서 17,576번의
클럭 내내 이 ``더하기 아니면 멈춤'' 규칙이 직접 뺄셈한 값과 같음을 확인한다.

\paragraph{릴레이 한 자리.} 0--25를 세는 릴레이 계수기 한 자리를 보자. 이진 계수기에서 1을 더하면 ``아래에서부터
연달아 1인 비트와 그 위의 첫 0비트''가 뒤집힌다(0111 + 1 = 1000). 뒤집힐 비트를 정하는 것을 \term{올림 사슬}{carry
chain}이라 한다. 문제는 릴레이가 뒤집히는 동안 그 비트를 보고 있는 올림 사슬도 흔들린다는 것이다. 뒤집힌 비트가
다시 사슬을 바꾸면 결과가 릴레이들의 빠르기 경쟁에 달린다.

\General 디지털 회로는 이 문제를 \term{마스터--슬레이브}{master--slave} 플립플롭으로 푼다. 기억 소자를 둘 두고
클럭의 두 상에서 번갈아 문을 연다. 한 상에서 앞 소자가 다음 값을 받아 두고(뒷 소자는 닫혀 있다), 다른
상에서 뒷 소자가 그것을 받는다(앞 소자는 닫혀 있다). 두 문이 동시에 열리지 않으므로 값이 한 클럭에 두 번
넘어가지 않는다. \ref{ch:sync}장의 ``클럭 에지에서만 바뀐다''는 성질을 이렇게 만든다.

\Ours 계수기 보드(\code{hardware/relay/counter-board})의 한 자리는 비트마다 릴레이 둘과 자리 전체에 하나, 모두
11개다(그림~\ref{fig:r-counter}).
\begin{itemize}
\item \textbf{슬레이브} KS$_i$: 값을 갖는다. 극 1이 +24~V를 \code{Q}$_i$(= 시프터의 제어선)에 잇고, 극 2가 자기
유지를 한다.
\item \textbf{마스터} KM$_i$: 평소에 늘 슬레이브의 \emph{반대} 값을 따라간다. ``이 비트를 뒤집으면 될 값''을
미리 계산해 두는 것이다. 마스터의 NC 접점들이 직렬로 올림 사슬을 이룬다. 마스터가 0이면(= 슬레이브가 1이면)
사슬이 다음 비트로 이어진다.
\item \textbf{감김 릴레이} KW: 값이 25(= 11001$_2$)가 아니면 켜져 있다. 코일은 비트 0, 3, 4의 반대값을 다이오드
OR로 받는다.
\end{itemize}
클럭은 두 상이다. 평소(\code{NCK})에는 마스터가 슬레이브를 거울처럼 따라가고 슬레이브는 자기 유지로 값을
붙든다. 확정 펄스(\code{CK}) 동안에는 사슬이 닿은 비트의 슬레이브만 마스터의 값을 받고, 그 비트의 마스터는
스스로를 붙들어 움직이지 않는다. 사슬이 \emph{클럭 동안 움직이지 않는} 마스터의 접점으로만 이루어져 있으므로
슬레이브가 바뀌어도 사슬은 흔들리지 않는다. 경쟁이 없다.

\begin{figure}[ht]
\centering
\begin{tikzpicture}[x=1cm,y=1cm,
    ks/.style={rcoil,minimum width=13mm,minimum height=7mm,font=\footnotesize},
    km/.style={rcoil,fill=cbluebg,minimum width=13mm,minimum height=7mm,font=\footnotesize}]
  \foreach \i in {0,1,2,3,4} {
    \node[ks] (S\i) at (2.2*\i,1.7) {KS$_\i$};
    \node[km] (M\i) at (2.2*\i,0) {KM$_\i$};
    \draw[->,cgreen] ([xshift=-2mm]S\i.south) -- node[left,font=\tiny,text=cgreen] {NOT} ([xshift=-2mm]M\i.north);
    \draw[->,corange] ([xshift=2mm]M\i.north) -- node[right,font=\tiny,text=corange] {CK} ([xshift=2mm]S\i.south);
    \draw[->] (S\i.north) -- ++(0,0.45) node[above,font=\scriptsize] {\code{Q}$_\i$};
  }
  \node[font=\footnotesize,anchor=east] (ain) at (-1.0,-1.3) {\code{AIN}};
  \draw (ain.east) -- (-0.7,-1.3);
  \foreach \i in {0,1,2,3,4} {
    \pgfmathtruncatemacro{\j}{\i+1}
    \draw (2.2*\i-0.7,-1.3) -- (2.2*\i-0.45,-1.3);
    \draw (2.2*\i-0.45,-1.3) -- (2.2*\i+0.25,-1.05);
    \draw (2.2*\i+0.25,-1.3) -- (2.2*\i+1.7,-1.3);
    \fill (2.2*\i-0.45,-1.3) circle (1pt); \fill (2.2*\i+0.25,-1.3) circle (1pt);
    \draw[dotted] (M\i.south) -- (2.2*\i-0.1,-1.15);
    \node[font=\tiny,anchor=north] at (2.2*\i-0.1,-1.35) {$A_\i$};
  }
  \node[font=\footnotesize,anchor=west] at (10.9,-1.3) {$A_5$ = 올림 \code{AOUT}};
  \node[note,anchor=north,align=center] at (4.4,-1.8) {올림 사슬: 마스터 $i$의 NC 접점. 마스터 $i$ = 0(슬레이브 $i$ = 1)이면 닫혀 $A_{i+1}$이 선다.\\
    $A_i$가 선 비트만 \code{CK} 동안 슬레이브가 마스터의 값(= 옛 값의 반대)을 받는다. 곧 그 비트가 뒤집힌다.};
  \node[ks,fill=credbg] (kw) at (11.7,0.85) {KW};
  \node[note,anchor=north,align=center] at (11.7,0.35) {값 $\ne 25$이면 켜짐.\\25에서 꺼져 비트 1, 2의\\거울을 끊는다};
\end{tikzpicture}
\caption{릴레이 계수기의 한 자리. 노란 상자가 값을 갖는 슬레이브, 파란 상자가 ``반대값''을 미리 계산해 두는
마스터다. 초록 화살표(평소, \code{NCK})는 마스터가 슬레이브의 반대를 따라가는 길, 주황 화살표(\code{CK})는
슬레이브가 마스터 값을 받는 길이다.}
\label{fig:r-counter}
\end{figure}

\paragraph{25 다음은 0.} 0--25를 세는 계수기는 25에서 0으로 돌아가야 한다. 보통은 26이 되는 순간을 감지해
지우지만, 그러면 아주 잠깐 26이라는 값이 시프터에 보인다. 이 설계는 감김 릴레이 KW 하나로 해결한다. 값이
25(11001$_2$)일 때만 KW가 꺼져 비트 1, 2의 마스터 거울 길을 끊는다. 그러면 평소라면 00110$_2$이었을 마스터가
00000이 된다. 마스터가 모두 0이므로 NC 사슬이 끝까지 닫혀 다섯 비트가 \emph{모두} 뒤집히고, 슬레이브는 모두
마스터 값 0을 받는다. 사슬의 끝 $A_5$가 서서 다음 자리로 올림을 넘긴다. 정상적인 값에서 마스터가 모두 0이
되려면 값이 31이어야 하므로 이 일은 25에서만 일어난다. 26은 한 번도 나타나지 않는다.

생성기에서 마스터와 슬레이브는 이렇게 적힌다. 접점 하나하나에 붙은 선 이름이 위의 설명과 같다.

\begin{lstlisting}[language=Python,caption={\code{gen\_counter.py}: 비트 $i$의 슬레이브와 마스터(\code{d}는 자리 번호).},label={lst:r-digit}]
        # slave: pole 1 = output pair Q/NQ from +24 V, pole 2 = rail Z -> (NC: mirror to master, NO: self-hold)
        relay(sh, path, f"K{d}0{i}", x, 60.0, f"Slave (state) relay bit {i}",
              dict(c1="+24V", no1=f"Q{i}", nc1=f"NQ{i}" if i in (0, 3, 4) else None, c2=z_net(i),
                   nc2=f"MIR{i}", no2=f"SH{i}", coil_p=f"SC{i}", coil_n="GND"))
        # master: pole 1 = carry chain A_i -> A_{i+1} (NC) and load/hold feed G_i (NO), pole 2 = Y_i = AIN & M_i
        relay(sh, path, f"K{d}1{i}", x, 135.0, f"Master relay bit {i} (= NOT slave)",
              dict(c1=a_net(i), nc1=a_net(i + 1), no1=f"G{i}", c2="AIN", nc2=None, no2=f"Y{i}",
                   coil_p=f"MC{i}", coil_n="GND"))
\end{lstlisting}

\paragraph{계수 규칙의 배치.} 자리들은 ``이번 클럭에 오르는가''를 \code{AIN}/\code{ZIN}이라는 상보 쌍의
선으로 받는다. 매 클럭 오르는 $l$의 \code{AIN}은 \code{CK} 그 자체이고, $m$의 \code{AIN}은 $l$ 자리의 올림이다.
$a_3$의 \code{AIN}은 ``\code{CK}이고 $l \ne 25$''인데, $l$ 자리의 마스터 접점이 이미 그 조건을 만들어 두므로
선 하나를 끌어오면 된다. 판정이 모두 클럭 동안 움직이지 않는 마스터에서 나오므로, 한 클럭의 모든 계수기가
같은 \code{CK} 펄스 안에서 동시에 움직이고 사이의 순서가 없다. \ref{ch:sync}장의 ``모든 레지스터가 같은
에지에서 동시에 바뀐다''를 릴레이로 만든 셈이다.

공통 보드가 $l$, $m$, $a_3$, $r$ 네 자리와 규칙을 맡고, 유닛 보드는 $a_2$, $a_1$ 두 자리를 받기만 한다.
유닛마다 다른 것은 프리셋($a_1 = k$, $a_2 = -k$)뿐이고, 그 값은 유닛마다 하나인 오프셋 다이얼이 준다. 이
다이얼의 동박 무늬를 이 기준 모델에서 직접 만드는 과정은 \ref{sec:gen-dial}절에서 본다.

\section{정리}

\begin{itemize}
\item 릴레이는 코일 전류로 접점을 옮기는 스위치다. 접점은 COM, NC(b접점), NO(a접점)이고, 이 프로젝트는 극이
둘인 전환 접점 릴레이(DPDT) 한 종류를 쓴다.
\item 접점 논리는 길 찾기다. 직렬은 AND, 병렬은 OR, NC는 NOT, 다이오드 여럿은 OR. 접점은 양방향이라 샛길에
주의한다.
\item 자기 유지 접점이 1비트 기억을 만든다. 끄려면 유지 길을 따로 끊는다.
\item 전환 접점은 먼저 끊고 나중에 잇는다. 정확히 하나를 골라야 하는 두 신호(\code{CK}/\code{NCK})는 한 접점에서
낸다.
\item 코일마다 역기전력 다이오드를 단다. 늦게 놓이는 대가가 있어, 빨리 놓여야 하면 제너를 더한다.
\item 접점 나무는 $n$글자를 $n-1$개의 전환 접점으로 고른다. 뿌리를 가장 낮은 비트에 두면 릴레이가 가장 적다.
\item 로터 = 밀기 + 고정 배선 + 되돌림. 밀기는 1·2·4·8·16칸 단의 직렬이고, 이웃한 밀기를 합쳐 유닛을
$S(a_1)W_RS(a_2)W_MS(a_3)W_LS(-l)$로 만든다.
\item 차이를 직접 세면 모든 계수기가 ``더하기 아니면 멈춤''만 한다. 릴레이 계수기는 마스터--슬레이브와 두 상의
클럭으로 경쟁 없이 센다.
\end{itemize}
